\documentclass[12pt]{exam}
\usepackage[margin=0.1in]{geometry}

%\documentclass[12pt]{exam}
\usepackage[utf8]{inputenc}
\usepackage{amsmath,amssymb}
\usepackage{graphicx}
\usepackage{enumitem}
\usepackage{multicol}
\usepackage{xcolor}
\usepackage{tikz}
\usetikzlibrary{arrows.meta,positioning,shapes.callouts}

\newcommand{\tfq}[3]{%
  \question[#1]\textbf{T/F:} #2%
  % Only show the choices when printing answers/keys
  \ifprintanswers
    \par\begin{oneparchoices}
      #3
    \end{oneparchoices}
  \fi
}

\makeatletter
\newcommand{\TFQ@choices}{} % temp holder

\newenvironment{tfqverb}[2]{%
  \def\TFQ@choices{#2}% store choices for use at end
  \question[#1]\textbf{T/F:}\ %
}{%
  \ifprintanswers
    \par\begin{oneparchoices}\TFQ@choices\end{oneparchoices}
  \fi
}
\makeatother

\printanswers
% optional styling for the correct choice:
\CorrectChoiceEmphasis{\bfseries}    % make the correct choice bold
% or e.g. \CorrectChoiceEmphasis{\color{red}\bfseries}


\newcommand{\SQLBlock}[1]{%
  \par\noindent
  \fbox{%
    \begin{minipage}[t]{0.97\linewidth}\ttfamily\small
      #1
    \end{minipage}%
  }%
}



\begin{document}
\section*{Review 2}

\begin{questions}
%---------------------------------TRUE/FALSE---------------------------------
\section*{Part I — True/False}

{
  \renewcommand{\choicelabel}{}
  \tfq{2}{A query with an aggregate function in the 'select' part of it and without a 'group by' statement returns only one tuple.}{\CorrectChoice True \choice False}
  \begin{tfqverb}{2}{\CorrectChoice True \choice False}
  Below query returns courses \((\mathit{course\_id}, \mathit{title}, \mathit{dept\_name}, \mathit{credits})\) that have more than one direct prerequisite.
  \begin{verbatim}
    select * 
    from course 
    where (select count(course_id) 
           from prereq 
           where prereq.course_id = course.course_id) >= 2;
  \end{verbatim}
  \end{tfqverb}
  \begin{tfqverb}{2}{\choice True \CorrectChoice False}
  Definition of the \textbf{prereq} table could be changed from
  \begin{verbatim}
    create table prereq
    (course_id varchar(8),
    prereq_id varchar(8),
    primary key (course_id, prereq_id),
    foreign key (course_id) references course(course_id) on delete cascade,
    foreign key (prereq_id) references course(course_id));
  \end{verbatim}
  to
  \begin{verbatim}
    create table prereq
    (course_id varchar(8),
    prereq_id varchar(8),
    primary key (course_id, prereq_id),
    foreign key (course_id) references course(course_id) on delete cascade,
    foreign key (prereq_id) references course(course_id) on delete cascade);
  \end{verbatim}
  \end{tfqverb}
  \begin{tfqverb}{2}{\CorrectChoice True \choice False}
  Consider the following scenario involving a view definition and an insertion operation. 
  Assume the insertion succeeds (no other instructor has id \texttt{'25566'} and a \texttt{Biology} department exists). 
  The count of rows returned by the last command remains unchanged before and after the insertion.
  \begin{verbatim}
  create view history_instructors as
  select * from instructor where dept_name = 'History';

  insert into history_instructors values
  ('25566', 'Brown', 'Biology', 100000);

  select count(*) from history_instructors;
  \end{verbatim}
  \end{tfqverb}
  \tfq{2}{Deletion of a tuple from the \textbf{department} relation may result in the deletion of some tuples from other relations.}{\choice True \CorrectChoice False}
  \tfq{2}{Once a transaction has executed commit work, its effects can no longer be undone by rollback work.}{\CorrectChoice True \choice False}
  \begin{tfqverb}{2}{\CorrectChoice True \choice False}
  In the result of the provided query, there might be instructors whose information \((\mathit{ID}, \mathit{name}, \mathit{dept\_name}, \mathit{salary})\) is not included.
  \begin{verbatim}
  select * from instructor natural left outer join teaches;
  \end{verbatim}
  \end{tfqverb}
  \begin{tfqverb}{2}{\choice True \CorrectChoice False}
  Below query can return a tuple with value 0 for the attribute \(\mathit{classroom\_no}\).
  \begin{verbatim}
  select building, count(distinct room_number) as classroom_no
  from classroom
  group by building;
  \end{verbatim}
  \end{tfqverb}

}


%---------------------------------MC---------------------------------
\section*{Part II — Multiple choice (choose one)}
\question[2] Assume that there is only one tuple in the \textbf{section} relation and no tuples in both \textbf{takes} and \textbf{student} relations. What is the minimal number of tuples that have to be added to the database so that the query below returns 2 tuples?

\begin{verbatim}
select * from student natural join takes;
\end{verbatim}

\begin{choices}
  \choice 2
  \choice 3
  \CorrectChoice 4
  \choice 5
\end{choices}




%---------------------------------MS---------------------------------
\section*{Part III — Multiple select (select one or more)}
\question[4] Which SQL queries retrieve the IDs and titles of \textbf{Comp. Sci.} courses offered in \textbf{Spring 2010}?
\begin{choices}

  \CorrectChoice
\begin{verbatim}
select course.course_id, title
from course, section
where course.course_id = section.course_id
  and course.dept_name = 'Comp. Sci.'
  and section.semester = 'Spring'
  and section.year = 2010;
\end{verbatim}

  \choice
\begin{verbatim}
select course.course_id, title
from course, section
where course.dept_name = 'Comp. Sci.'
  and section.semester = 'Spring'
  and section.year = 2010;
\end{verbatim}

  \choice
\begin{verbatim}
select course.course_id, title
from course, section, department
where course.course_id = section.course_id
  and department.dept_name = 'Comp. Sci.'
  and section.semester = 'Spring'
  and section.year = 2010;
\end{verbatim}

  \choice
\begin{verbatim}
select course_id, title
from course
where dept_name = 'Comp. Sci.'
  and semester = 'Spring'
  and year = 2010;
\end{verbatim}

\end{choices}

\question[5] Students (their \(\mathit{ID}, \mathit{name}, \mathit{dept\_name}, \mathit{tot\_cred}\)) who don't have advisors.
\begin{choices}

  \choice
\begin{verbatim}
with temp(ID) as
  (select s_id
   from advisor
   where i_id is null)
select *
from student natural join temp;
\end{verbatim}

  \choice
\begin{verbatim}
select ID, name, dept_name, tot_cred
from student natural join advisor
where i_id is null;
\end{verbatim}

  \CorrectChoice
\begin{verbatim}
select ID, name, dept_name, tot_cred
from student left outer join advisor
  on (student.ID = advisor.s_id)
where advisor.s_id is null;
\end{verbatim}

  \choice
\begin{verbatim}
select *
from student
where ID not in (select i_id from advisor);
\end{verbatim}

  \choice
\begin{verbatim}
select ID, name, dept_name, tot_cred
from student natural join advisor natural join instructor
where ID is null;
\end{verbatim}

\end{choices}

\question[4] Retrieve all course sections (identified by \((\mathit{course\_id}, \mathit{sec\_id}, \mathit{semester}, \mathit{year})\)) where the total enrollment exceeds the seating capacity of the assigned classroom. \emph{(Select all that apply.)}

\begin{choices}

  \CorrectChoice
\begin{verbatim}
select s.course_id, s.sec_id, s.semester, s.year
from section s
join takes t
  on s.course_id = t.course_id
 and s.sec_id    = t.sec_id
 and s.semester  = t.semester
 and s.year      = t.year
join classroom c
  on s.building = c.building
 and s.room_number = c.room_number
group by s.course_id, s.sec_id, s.semester, s.year, c.capacity
having count(t.ID) > c.capacity;
\end{verbatim}

  \CorrectChoice
\begin{verbatim}
with enroll(course_id, sec_id, semester, year, students_no) as (
  select course_id, sec_id, semester, year, count(*)
  from takes
  group by course_id, sec_id, semester, year
)
select s.course_id, s.sec_id, s.semester, s.year
from enroll e
join section s using (course_id, sec_id, semester, year)
join classroom c
  on s.building = c.building and s.room_number = c.room_number
where e.students_no > c.capacity;
\end{verbatim}

  \choice
\begin{verbatim}
with temp(course_id, sec_id, semester, year, students_no) as (
  select course_id, sec_id, semester, year, count(ID)
  from section natural join takes
  group by course_id, sec_id, semester, year
)
select course_id, sec_id, semester, year
from temp natural join classroom
where capacity < students_no;
\end{verbatim}

  \choice
\begin{verbatim}
select course_id, sec_id, semester, year
from section natural join takes
group by course_id, sec_id, semester, year
having count(*) > capacity;
\end{verbatim}

\end{choices}

\question[4] Assume the following roles/grants were executed on the University database server for users \texttt{Grey}, \texttt{Amit}, \texttt{Green}, and \texttt{Brown}. Which user(s) have sufficient privileges to move all \textbf{Computer Science} department students who enrolled in \texttt{CS-190} in \textbf{Spring 2009} to section \texttt{2} of the course (assume section 2 already exists)? In other words, which user(s) can successfully execute the query below?


\begin{verbatim}
update takes
set sec_id = 2
where id in (select id from student where dept_name = 'Comp. Sci.')
  and course_id = 'CS-190'
  and semester = 'Spring'
  and year = 2009;
\end{verbatim}

Executed Commands:
\begin{verbatim}
create role 'teaching_assistant'@'localhost';
create role 'lecturer'@'localhost';
create role 'dean'@'localhost';

grant select on university.takes to 'teaching_assistant'@'localhost';
grant update on university.takes to 'teaching_assistant'@'localhost';

grant select on university.student to 'lecturer'@'localhost';

grant select on university.instructor to 'dean'@'localhost';
grant 'teaching_assistant'@'localhost' to 'lecturer'@'localhost';
grant 'lecturer'@'localhost'        to 'dean'@'localhost';

grant 'dean'@'localhost' to 'Grey'@'localhost';
grant 'dean'@'localhost' to 'Amit'@'localhost';
grant update on university.instructor to 'Amit'@'localhost';

grant update on university.student to 'Brown'@'localhost';
grant 'lecturer'@'localhost' to 'Brown'@'localhost';

grant 'teaching_assistant'@'localhost' to 'Green'@'localhost';
\end{verbatim}

Hint: To test this scenario in MySQL, make sure to enable the default roles for each user by executing:
\begin{verbatim}
SET DEFAULT ROLE ALL TO 'Brown'@'localhost';
SET DEFAULT ROLE ALL TO 'Grey'@'localhost';
SET DEFAULT ROLE ALL TO 'Amit'@'localhost';
SET DEFAULT ROLE ALL TO 'Green'@'localhost';
\end{verbatim}

Important: To correctly test this scenario, each user must run the following two commands&nbsp;after logging in, before executing the UPDATE query:

1. Activate all granted roles for the current session:
SET ROLE ALL;

2. Select the appropriate database:
USE university;

\begin{choices}
  \CorrectChoice \texttt{Grey}
  \CorrectChoice \texttt{Amit}
  \choice        \texttt{Green}
  \CorrectChoice \texttt{Brown}
\end{choices}





\end{questions}


\end{document}
